%%
%%
%% anx_latex.tex for  in cours/outils_gnu
%%
%% Made by Philippe THIERRY
%% Login   <Philippe THIERRYreseau-libre.net>
%%
%% Started on  sam. 18 sept. 2010 11:43:34 CEST Philippe THIERRY
%% Last update sam. 18 sept. 2010 11:50:40 CEST Philippe THIERRY
%%

\chapter{Exemple de documentation logicielle r�dig�e en \LaTeX}

\paragraph{}
Le script \ref{lst:latex_example} est un exemple (r�el) de documentation \LaTeX d'un projet Open-Source.

\begin{lstlisting}[caption={Documentation logicielle type r�dig�e en \LaTeX},label=lst:latex_example]
%%
%%
%% requirements.tex for libcommon in libcommon/doc/latex/includes
%%
%% Made by phil
%% Login   <philreseau-libre.net>
%%
%% Started on  mer. 09 sept. 2009 23:05:40 CEST phil
%% Last update mer. 02 juin 2010 19:58:56 CEST Philippe THIERRY
%%
\documentclass[a4paper]{book}
\usepackage{a4wide}
\usepackage[Lenny]{fncychap}
\usepackage{makeidx}
\usepackage{fancyhdr}
\usepackage[dvips]{graphicx}
\usepackage{multirow}
\usepackage{multicol}
\usepackage{float}
\usepackage{textcomp}
\usepackage{alltt}
\usepackage{times}
\usepackage{ifpdf}
\usepackage{lastpage}
%\usepackage{calc}
\usepackage{datetime}
\usepackage{chngpage}
\usepackage{tabularx}
\usepackage{listings}

\usepackage{styles/macros}

\usepackage{styles/requirement}
\usepackage[toc,page]{appendix}
\usepackage{color}
\usepackage[hundred]{vrsion}
\ifpdf
\usepackage[pdftex,
            pagebackref=true,
            colorlinks=true,
            linkcolor=blue
           ]{hyperref}
\else
\usepackage[ps2pdf,
            pagebackref=true,
            colorlinks=true,
            linkcolor=blue
           ]{hyperref}
\usepackage{pspicture}
\fi

\usepackage{titlesec}

% charts support
\usepackage{pdftricks}
\begin{psinputs}
  \usepackage{pstricks}
  \usepackage{multido}
  \usepackage{epsfig}
  \usepackage{pst-grad} % For gradients
  \usepackage{pst-plot} % For axes
\end{psinputs}

% Add a background picture to first page
\usepackage{eso-pic}
% coloration of section and subsection titiles
\definecolor{seccol}{rgb}{0.8,0,0}
\definecolor{subseccol}{rgb}{0,0,0.7}

\titleformat{\section}
{\color{seccol}\normalfont\Large\bfseries}{\thesection}{1em}{}

\titleformat{\subsection}
{\color{subseccol}\normalfont\large\bfseries}{\thesubsection}{1em}{}

\makeindex
\setcounter{tocdepth}{1}
\renewcommand{\footrulewidth}{0.4pt}

% add code support
\definecolor{colKeys}{rgb}{0,0,1}
\definecolor{colIdentifier}{rgb}{0,0,0}
\definecolor{colComments}{rgb}{0,0.5,1}
\definecolor{colString}{rgb}{0.6,0.1,0.1}
\definecolor{colList}{rgb}{0.9,0.9,0.9}

%%%% debut macro %%%%
\newenvironment{changemargin}[2]{\begin{list}{}{%
\addtolength{\leftmargin}{#1}%
\addtolength{\rightmargin}{#2}%
}\item }{\end{list}}
%%%% fin macro %%%%

\renewcommand{\chaptermark}[1]{%
\markboth{\MakeUppercase{%
\thechapter. #1}}{}}

\definecolor{mail}{rgb}{1,0,0}

% update table of contents depth
\setcounter{tocdepth}{3} % Dans la table des matieres
\setcounter{secnumdepth}{3} % Avec un numero.

\pagestyle{fancy}
\fancyhf{}

%%
%%
%% titlepage.tex for doctorat in /home/phil/Travail/Scolarité/Doctorat/doc/tex
%%
%% Made by Philippe THIERRY
%% Login   <Philippe THIERRYreseau-libre.net>
%%
%% Started on  mar. 24 nov. 2009 19:50:24 CET Philippe THIERRY
%% Last update Fri Sep 10 14:57:58 2010 Philippe THIERRY
%%

\newcommand{\thesisTitle}{
\changepage{3cm}{0cm}{0cm}{0cm}{0cm}{0cm}{-1cm}{0cm}{0cm}
\begin{center}
\begin{adjustwidth}{-1cm}{-1cm}
\begin{tabular}{p{9cm}p{8.5cm}}
\begin{flushleft}\includegraphics[width=5cm]{pictures/ECE_LOGO_FR}\end{flushleft} &
%\begin{flushright}\includegraphics[width=5.5cm]{pictures/thales}\end{flushright} 
\\
\end{tabular}\\
\end{adjustwidth}
%\huge{\textbf{ECOLE CENTRALE D'ELECTRONIQUE\\
%PARIS}}\\
\vspace{1.1cm}
\large{\textsc{}}\\
\vspace{1.1cm}
Document �crit par\\
\vspace{0.5cm}
\Large{\textbf{Philippe THIERRY}}\\
\vspace{0.8cm}
dans le cadre des projets de l'ECE\\
\vspace{0.5cm}
\large{\textsc{}}\\
\vspace{1.8cm}
Sujet\\
\vspace{1.0cm}
\begin{adjustwidth}{-1cm}{-1cm}
\begin{tabular}[width=19cm]{||c||}
\hline
\hline
 \\
\huge{\textsc{Les outils GNU pour la production}}\\
\huge{\textsc{Formation initiale aux outils de production et}}\\
\huge{\textsc{�tudes de cas r�els}}\\
  \\
\hline
\hline
\end{tabular}\\
\end{adjustwidth}
\vspace{2.5cm}
\small{\it version 1.0}\\
\vspace{1cm}
\normalsize{Avec la participation de}\\
\vspace{0.7cm}
\begin{tabular}{lll}
\large{Jean-Marc} & \large{LACROIX} & \large{jeanmarc.lacroix@free.fr} \\
\large{Sylvain} & \large{LEROY} & \large{sylvain@unmondelibre.fr} \\
\end{tabular}
\end{center}
\newpage
\changepage{-3cm}{0cm}{0cm}{0cm}{0cm}{0cm}{1cm}{0cm}{0cm}
}


\begin{document}
% Declare ref counter. This counter is used to increment the [ REF X ]
% values, for successive elements defined as a reference
\newcounter{ref}
\stepcounter{ref}

\renewcommand{\headrulewidth}{0pt}
\renewcommand{\footrulewidth}{0pt}

% Generate the titlepage background using given picture
\AddToShipoutPicture*{\BackgroundPic{share/img/libcommon-background.png}}
\libcommontitle{the requirements cookbook}{0.98.8}

\cleardoublepage
\renewcommand{\headrulewidth}{0.5pt}
\fancyhf{}
\pagestyle{fancy}
\renewcommand{\footrulewidth}{1pt}
\fancyhead[LO]{\uppercase{\leftmark}}
\fancyhead[RO]{\rightmark}
\fancyhead[LE]{\leftmark}
\fancyhead[RE]{\rightmark}
\fancyfoot[LE]{\thepage/\pageref{LastPage}}
\fancyfoot[RO]{\thepage/\pageref{LastPage}}
\fancyfoot[C]{libcommon v.0.98.8 - generated on September, 15 2010}
\pagenumbering{roman}
\tableofcontents
\cleardoublepage
\listoffigures
\cleardoublepage
\pagenumbering{arabic}

\begin{minipage}{0.8\linewidth}
\vspace{2cm}
{\Huge Abstract\\}
\vspace{2cm}
{
\paragraph{}
\it
This document describes all the input requirements of the libcommon project.\\
This document is in compliance with the DO178B-DAL C, presenting the system
requirements\footnote{see The DO-178B description, part {\it system
requirements}}. This document is the concatenation of all the specialized
design \& requirements documents.
}
\end{minipage}
\newpage

%%
%%
%% introduction.tex for  in /doctorat/ece/partenariat/cours/outils_gnu
%%
%% Made by Philippe THIERRY
%% Login   <Philippe THIERRYreseau-libre.net>
%%
%% Started on  Wed Sep  1 18:07:05 2010 Philippe THIERRY
%% Last update mer. 01 sept. 2010 19:52:38 CEST Philippe THIERRY
%%

\chapter{Introduction}

\paragraph{}
Le but de ce document et de fournir une base de connaissance suffisante permettant d'apréhender l'architecture des
éléments logiciels utilisé dans les environnements UNIX/Linux.

\paragraph{}
Le nombre de logiciels s'appuyant sur le système de production décrit ci-après est très important, et l'est de plus
en plus avec l'arrivée des nouvelles plateforme logicielles embarquées dans la téléphonie ou encore la télévision,
s'appuyant pour la plupart sur un socle GNU/Linux.\\
On retrouve également ces mêmes logiciels dans les systèmes réseaux, comme les {\it box}, ou encore les routeurs et
appliances (hors CISCO IOS).

\paragraph{}
Dans le métier des télécommunications ou encore chez les industriels de la téléphonie ou les fournisseurs de
logiciels embarqués, les projets se basent de plus en plus fréquement sur le système de production standardisé par
les UNIX, sous sa version libre GNU (Gnu Compiler Collection). Ce document en décrit les arcanes et met en lumière
les éléments de complexité afin de fournir une base nécessaire à l'étude des grands logiciels.


% code quality
\input{code_quality.tex}
% software architecture
\input{software_architecture.tex}
% software production
\input{software_production.tex}
% code quality
\input{software_test.tex}
% software documentation
\input{software_documentation.tex}

\printindex
\end{document}
\end{lstlisting}
